\section{Related Works}
\label{sec:related_works}

In this section, we provide a comprehensive review of prior studies in layout analysis, text recognition, and the long-tailed distribution problem, with a particular emphasis on their relevance to historical text recognition tasks. Various approaches, including both classical optical character recognition (OCR) systems and more recent deep learning-based techniques, have been explored to address the unique challenges posed by historical documents~\cite{liu2019survey, bishop2006pattern}. Nevertheless, the intricacies involved in digitizing centuries-old texts, which often include noise, uneven ink spread, and archaic writing styles, have propelled researchers to devise specialized strategies. This section outlines the challenges inherent in current methods and sets the stage for our proposed solution, while emphasizing that effective historical text recognition necessitates a careful consideration of layout complexities, language-specific characteristics, and the long-tailed distribution of characters~\cite{jla}.

\subsection{Historical Text Recognition}

Document digitization plays an essential role in preserving cultural heritage by converting physical documents into digital formats, enabling efficient storage, retrieval, and analysis. The transformation of historical texts into machine-readable data has been significantly facilitated by OCR technologies, which range from traditional rule-based methods to advanced deep neural networks~\cite{conway2010, yadav2021handwritten}. Among its key components, text recognition has garnered significant attention as it is central to transforming digitized documents into formats that allow historians, linguists, and archivists to conduct comprehensive research~\cite{jla}.

Historical text recognition methods are generally categorized into two main approaches: character-based methods and sequence-based methods. Both offer unique strengths but also face challenges when applied to historical documents, including the presence of rare symbols, blurred ink traces, and nonlinear text flow~\cite{chetarly22}. These difficulties underscore the need for robust and adaptable approaches that can handle the variations introduced by time and usage.

\subsubsection{Character-Based Methods}

Character-based methods focus on processing individual characters by decomposing the recognition pipeline into locating characters, recognizing them, and grouping them into meaningful text lines~\cite{papytwin}. This strategy has been widely studied due to its ability to explicitly model character-level details, which can be advantageous when dealing with highly complex or irregular character shapes. However, it remains particularly challenging in historical contexts because the documents often exhibit degraded quality with broken or smudged characters caused by age-related wear and imperfect printing processes~\cite{broken}. Historical materials may also include handwritten texts with widely varying writing styles, introducing a layer of complexity as the differences in stroke density and curvature can hinder recognition accuracy~\cite{obc306}. Furthermore, the frequency of characters is often heavily skewed, implying that a few commonly used characters dominate the distribution while a vast number of rare characters rarely appear~\cite{fewran}. These rare characters may include novel symbols or archaic ligatures~\cite{hde,ligarature} that never appear in standard training datasets, thus complicating the development of universal or language-agnostic models. Despite these difficulties, the character-based approach remains a foundational method because its modularity allows for fine-grained control over each step of the recognition pipeline and provides a mechanism to explicitly address rare or newly discovered characters during the digitization of historically significant texts.

\subsubsection{Sequence-Based Methods}

Sequence-based methods aim to recognize entire lines of text directly from images, bypassing the need for explicit character segmentation~\cite{eccvfork,jinic21}. In these methods, a network is typically trained with line-level annotations, which are easier to obtain in large quantities compared to the detailed labeling required for character-level training~\cite{adams2019sequence}. Although sequence-based methods have reduced annotation burdens and demonstrated remarkable performance in modern OCR tasks, they still face limitations. First, they struggle to handle rare characters efficiently because imbalanced datasets often overemphasize the few high-frequency symbols while neglecting many low-frequency ones. Second, they exhibit sensitivity to variations in text layout and structure, which is especially problematic in historical documents with intricate page designs or non-standard formatting. Third, their reliance on implicit representations makes it difficult to generalize to unseen characters or to explicitly model novel symbols, thereby limiting their usefulness when encountering entirely new character shapes. While sequence-based methods are appealing for large-scale digitization efforts that seek to rapidly process line-level data, tasks that require detailed character-level analysis can suffer from the reduced interpretability and adaptability of these models. In this work, we adopt a character-oriented perspective to address the unique challenges of historical text recognition because it allows more direct handling of long-tailed and novel character distributions.

\subsection{The Long-Tailed Distribution Problem}

The long-tailed distribution problem is a pervasive issue in machine learning, particularly in tasks where a few classes dominate the dataset (referred to as head classes), while a majority of classes are underrepresented (referred to as tail classes)~\cite{tailsurvey}. This imbalance has far-reaching implications for model training, making it difficult to achieve high accuracy for underrepresented classes. Standard optimization methods often overfit to the head classes, capturing only limited nuances from rare classes. In text recognition, this skewed distribution leads to misclassification of rare characters, an issue that is magnified in historical text corpora where unique or archaic symbols appear infrequently~\cite{garcia2018improved}.

\subsubsection{Existing Solutions}

Existing solutions to the long-tailed problem can be situated within three primary strategies that seek to balance the representation of head and tail classes, thereby improving the overall performance of recognition systems. Data-side approaches focus on balancing the data distribution by introducing class-balancing sampling, which allocates more training emphasis on underrepresented characters~\cite{upsam}, or by adopting data augmentation techniques, such as generating synthetic samples to bolster the diversity of tail classes~\cite{cutmix, smote, hzsl}. Other strategies involve optimization-side approaches, where loss reweighting schemes allocate higher importance to rare classes~\cite{focal,otem} and additional regularization terms guide the model toward learning more generalizable features~\cite{fudanvi,sanicdar23,logvar}. Finally, model-side approaches concentrate on architectural modifications, including ensemble methods that aggregate predictions from multiple specialized models~\cite{fle,flmoe}, and classifier adaptations that adjust decision boundaries to address imbalanced distributions~\cite{normrob}. Each line of research offers distinct advantages, but combining these perspectives often leads to more robust solutions. Indeed, recent progress suggests that hybrid approaches that integrate data augmentation with reweighted loss functions or specialized network modules can achieve state-of-the-art results for long-tailed recognition tasks~\cite{chen20dgan}.

\subsubsection{Zero-Shot Learning and Its Relevance}

Zero-shot learning presents an extreme case of the long-tailed problem, where certain classes in the test set are absent from the training phase entirely~\cite{gzsl-survey,olt}. By leveraging shared attributes between head and tail classes, zero-shot methods enable models to generalize to unseen categories without direct sample supervision~\cite{shen22}. This capability is particularly valuable in historical text recognition contexts, where new or archaic characters may suddenly appear in a document. When used judiciously, zero-shot learning frameworks reduce the reliance on meticulously annotated data, and they can exploit the compositional nature of characters, such as the shared radicals or strokes in handwritten scripts. Despite the promise of these techniques, they also require careful architectural design to ensure the accurate transfer of learned knowledge to genuinely new symbols.

\subsection{Long Tails in Historical Text Recognition}

Historical text recognition presents unique challenges in addressing long-tailed distributions. The frequency distribution of characters in historical corpora often reflects natural language use, where common characters dominate the text while rare symbols occur sporadically~\cite{yang2022survey}. This phenomenon is exacerbated by the limited availability of annotated historical datasets, which not only restricts the variety of character exemplars but also intensifies the scarcity of training samples for tail classes. Researchers have attempted to alleviate these difficulties by exploiting shared knowledge between head and tail classes, such as radical composition~\cite{obcmk2,sanicdar23,fudanvi} or stroke-based features that capture fundamental structural elements of characters~\cite{taktak}. Although these approaches demonstrate improvements, they typically rely on detailed annotations that demand expert knowledge and are not easily transferable across different scripts or languages.

To overcome these limitations, we propose a novel approach that implicitly leverages the shared components of characters within a carefully designed backbone network. By emphasizing part-level features during feature extraction, our method reduces reliance on explicit annotations while enhancing performance for tail classes and novel characters. In particular, by focusing on the compositional structure of characters rather than their overall shape, we improve the model's ability to generalize to unusual or unseen symbols that often arise in historical documents. This approach addresses the core challenges of historical text recognition in a more general and scalable manner, offering a pathway to digitize vast collections of historical texts that exhibit highly imbalanced and evolving character distributions.

%-------------------------------
% Additional References in BibTeX Format
%-------------------------------
% Please insert the following references into your .bib file or reference list as needed.
